\documentclass[10pt,a4paper]{amsart}
\usepackage[margin=19mm]{geometry}
\usepackage{amsmath,amssymb,mathtools}
\usepackage[scr=rsfs,cal=euler]{mathalfa}
\usepackage{xfrac}
\usepackage[unicode,hidelinks,bookmarks=false]{hyperref}
\newcommand{\ie}{i.e.\ }
\newcommand{\cf}{cf.\ }
\newcommand{\resp}{respectively\ }
\newcommand{\Subsection}{\subsection}
\providecommand{\nobreakdash}{\nobreak\hbox{-}}
\setlength{\emergencystretch}{2em}
\multlinegap0pt
\usepackage{bm}
\usepackage{bbm}
\usepackage{dsfont}
\usepackage{xspace}
\usepackage{cancel}
\usepackage{mathtools}
\usepackage{amsthm}
\makeatletter
\@ifundefined{thm}{\newtheorem{thm}{Theorem}}{}
\@ifundefined{lem}{\newtheorem{lem}[thm]{Lemma}}{}
\makeatother
\newcommand{\Bb}{\overline{\mathcal{B}}}
\newcommand{\Hb}{\overline{\mathcal{H}}}
\newcommand{\Mb}{\overline{\M}}
\newcommand{\M}{\mathcal{M}}
\newcommand{\RH}{\mathsf{RH}}
\renewcommand{\H}{\mathsf{H}}
\renewcommand{\tilde}{\widetilde}
\title[The tautological ring of $\overline{\mathcal{M}}\_{g,n}$ is r]{The tautological ring of $\overline{\mathcal{M}}\_{g,n}$ is rarely Gorenstein (Lemma 19)}
\author{Samir Canning}
\date{Source: arXiv:2406.10516v1 (CC BY 4.0)}
\begin{document}
\maketitle

\noindent\emph{Source and licence.} The tautological ring of $\overline{\mathcal{M}}_{g,n}$ is rarely Gorenstein. Version arXiv:2406.10516v1 (CC BY 4.0), distributed under the Creative Commons Attribution 4.0 International licence, as recorded in the arXiv licence metadata for this exact version and shown on the abstract page https://arxiv.org/abs/2406.10516v1. Full rights notice: licenses/arxiv-2406.10516-CC-BY-4.0.txt.\par\smallskip

\noindent\textit{Editorial scope.} Counted unit: the Lem printed as Lemma 19 of the article (source0.tex lines 528--542, source label \texttt{tautterms}), delivered together with the article's own complete original proof of it (source0.tex lines 543--567, one \texttt{proof} environment of 25 lines). No mathematics is written, repaired or re-derived: statement and proof are the article's own text line for line. Required context retained uncounted: hyperelliptictaut (lines 355--357); PetersenTommasi (lines 363--370). Every definition, display and label of those lines is retained unchanged. Omitted from this package and not redistributed: 1--354, 358--362, 371--527, 568--634. Nothing inside a retained excerpt is omitted or altered: every retained body is delivered exactly as the article's lines are written, so no omission receipt of any kind accompanies this package. Citations: the retained excerpts carry no citation command at all, so no bibliography entry of the article is transcribed and no reference is redistributed. Reference adaptation: none was needed and none was made; every reference inside the delivered unit resolves inside it, so no unresolved reference or citation is produced. Printed slips kept literal: no printed word, symbol, hypothesis or number of the counted statement or of its proof was altered. Layout and class: the article's own class is \texttt{amsart} with packages including tikz-cd, enumitem, amssymb, amsmath, amscd, mathrsfs, url, cite, fullpage, hyperref, verbatim, comment, fancyvrb, fvextra; the wrapper is the lane's pinned \texttt{amsart} editorial wrapper at 10pt on A4, which restates the theorem-like environments the retained text needs and adds the article's own macro definitions that the retained text needs (\texttt{\textbackslash Bb}, \texttt{\textbackslash H}, \texttt{\textbackslash Hb}, \texttt{\textbackslash M}, \texttt{\textbackslash Mb}, \texttt{\textbackslash RH}, \texttt{\textbackslash tilde}), with extra wrapper packages (bm, bbm, dsfont, xspace, cancel, mathtools). The delivered numbering is the unit's own. Assets: no retained span contains a figure environment, a graphics-inclusion command, a PDF-page inclusion, a table or a TikZ picture; no image file is copied into this package, the package declares no asset dependency and no image file is redistributed with it. Licence: version arXiv:2406.10516v1 of this article is distributed under the Creative Commons Attribution 4.0 International licence, as recorded in the arXiv licence metadata for this exact version and shown on the abstract page https://arxiv.org/abs/2406.10516v1. A full rights notice is delivered at licenses/arxiv-2406.10516-CC-BY-4.0.txt. Characters: every retained excerpt is pure ASCII, so the delivered engine, XeLaTeX with its default Latin Modern text font, reports no missing character.
\begin{thm}[Faber--Pandharipande]\label{hyperelliptictaut}
The class $[\Bb_{g\rightarrow 0, n_1,2n_2}]$ is tautological for any $g,n_1,n_2$.
\end{thm}

\begin{thm}[Petersen--Tommasi, Petersen]\label{PetersenTommasi}
Let $k\in \mathbb{N}$ be even.
    \begin{enumerate}
        \item For $n<20$, $\RH^{k}(\Mb_{2,n})=\H^{k}(\Mb_{2,n})$.
        \item For $k\neq 22$, $\RH^{k}(\Mb_{2,20})=\H^{k}(\Mb_{2,20})$.
        \item The image of the pushforward map $\H^{k-2}(\tilde{\partial\M_{2,20}})\rightarrow \H^{k}(\Mb_{2,20})$ lies in $\RH^k(\Mb_{2,20})$, where $\tilde{\partial\M_{2,20}}$ is the normalization of the boundary.
    \end{enumerate}
\end{thm}

\begin{lem}\label{tautterms}
    Suppose $(\Gamma,G,f)\in \mathfrak{H}_B$ satisfies any of the following properties:

    \begin{enumerate}
        \item $\Hb_{(\Gamma,G)}$ is not of the expected codimension;
        \item $\Gamma$ has no vertices of genus $2$;
        \item $G$ acts nontrivially on the vertices;
        \item The quotient graph $\Gamma/G$ has first Betti number $1$;
        \item $\Gamma$ has loop edges;
        \item $G$ acts trivially on an edge;
        \item $\Gamma$ has vertices $v_1,v_2$ connected by $n\geq 1$ edges and $n\neq 2$.
    
    \end{enumerate}
Then $\pi_{B*}\left(c_{\mathrm{top}}(E_f)\cdot \phi_{f*}[\Hb_{(\Gamma,G)}]\right)$ is tautological.
\end{lem}

\begin{proof}
    Suppose $(\Gamma,G,f)$ is such that (1) holds. Then the excess bundle $E_f$ is of rank at least $1$. Because $c_{\mathrm{top}}(E_f)$ is pulled back from $\Mb_{A}$, the term
    \[
    \pi_{B*}\left(c_{\mathrm{top}}(E_f)\cdot \phi_{f*}[\Hb_{(\Gamma,G)}]\right)
    \]
    is a product of algebraic classes of lower codimension on $\Mb_{A}=\Mb_{2,20}$ by the projection formula. Hence, by Theorem \ref{PetersenTommasi}(2), $\pi_{B*}\left(c_{\mathrm{top}}(E_f)\cdot \phi_{f*}[\Hb_{(\Gamma,G)}]\right)$ is tautological. 
    
    From now on, we suppose $(\Gamma,G,f)$ does not satisfy property (1). Then there is no excess contribution, and we study only the terms of the form $\pi_{B*}\phi_{f*}[\Hb_{(\Gamma,G)}]$. Now suppose $(\Gamma,G,f)$ is such that (2) holds. First, assume that $\Gamma$ has a vertex of genus $1$. Then $\pi_{B*}\phi_{f*}[\Hb_{(\Gamma,G)}]$ is supported on the boundary of $\Mb_{A}$, and hence is tautological by Theorem \ref{PetersenTommasi}(3). We argue similarly if all vertices of $\Gamma$ are of genus $0$.

    From now on, we assume $(\Gamma,G,f)$ does not satisfy property (2), and so has a unique vertex of genus $2$. Suppose that $(\Gamma,G,f)$ is such that $G$ acts non-trivially on the vertices of $\Gamma$. We pick a pair $v_1$ and $v_2$ of genus $0$ vertices interchanged by the $G$ action. The $G$ action fixes all of the legs because the monodromy data is $\xi=(1^{22})$. Therefore, there can be no legs attached to $v_1$ and $v_2$. Under the pushforward $\pi_{B*}$ forgetting the legs, the vertices $v_1$ and $v_2$ are not contracted, meaning $\pi_{B*}\phi_{f*}[\Hb_{(\Gamma,G)}]$ is supported on the boundary of $\Mb_{A}$, and is thus tautological.

    From now on, we assume $(\Gamma,G,f)$ does not satisfy property (3). Suppose that $\Gamma/G$ has first Betti number $1$. It follows that every vertex of the quotient graph is of genus $0$, as the target curve has genus $1$. Therefore, in the product decomposition 
    \[
    \Hb_{(\Gamma,G)}=\prod_{v\in V} \Hb_{g(v),G_v,\xi_v},
    \]
    the term corresponding to the genus $2$ vertex parametrizes admissible hyperelliptic covers. Such spaces contribute only tautological classes by Theorem \ref{hyperelliptictaut}. 
    
    From now on, we assume $(\Gamma,G,f)$ does not satisfy property (4). Suppose that $\Gamma$ has a vertex with a loop edge. The quotient graph $\Gamma/G$ must have a loop edge because $G$ acts trivially on the vertices. This contradicts the assumption that $(\Gamma,G,f)$ does not satisfy property (4). 

    From now on, we assume $(\Gamma,G,f)$ does not satisfy property (5). Consider the case when $\Gamma$ has vertices $v_1$ and $v_2$ connected by an edge $e$ that is fixed under the $G$ action. Because the $B$ structure is generic, $e$ must be an edge of $B$. But $e$ is not a loop edge, whereas all the edges of $B$ are. Therefore, there must exist another path in $\Gamma$ connecting $v_1$ and $v_2$, in which all constituent edges are contracted by the $B$ structure $f:\Gamma\rightarrow B$. But then the quotient graph $\Gamma/G$ must have nonzero first Betti number, as the images of $v_1$ and $v_2$ are connected by at least two edges, contradicting that $(\Gamma,G,f)$ does not satisfy property (4).

    
    From now on, we assume $(\Gamma,G,f)$ does not satisfy property (6). Suppose that $\Gamma$ has vertices $v_1$ and $v_2$ connected by $n\geq 1$ edges $e_1,\dots,e_n$. If $n$ is odd, then because $G$ acts trivially on the vertices, it must also fix an edge, violating the assumption that it does not satisfy property (6). Therefore, $n$ is even. If $n\geq 4$, then the quotient graph $\Gamma/G$ has first Betti number at least $1$, violating the assumption that it does not satisfy property (4).

\end{proof}

\end{document}
